\problem{2.55}{Numerical shooting for the oscillator ground state}
With $q=\sqrt{m\omega/\hbar}\,x$ and
$K=2E/(\hbar\omega)$, the dimensionless equation is
\begin{equation}
u''(q)=(q^2-K)u(q).
\end{equation}
For an even ground state use $u(0)=1,u'(0)=0$,
integrate to a large cutoff $q=b$, and vary $K$
until the growing asymptotic component changes sign
and the tail becomes decaying. A robust alternative
is to match logarithmic derivatives of inward and
outward solutions, reducing sensitivity to the
exponentially growing solution.
The exact result provides a stringent five-digit check:
\begin{equation}
\boxed{K_0=1.0000,\qquad
E_0=0.50000\,\hbar\omega}.
\end{equation}
\check{A plot that appears to decay on a finite interval
does not by itself prove an eigenvalue. Verify node
count and cutoff convergence.}

\problem{2.56}{First three excited oscillator energies by shooting}
Solve the same equation as Problem 2.55.
For odd parity impose $u(0)=0,u'(0)=1$;
for even parity impose $u(0)=1,u'(0)=0$.
In increasing order the first three excited
solutions have one, two, and three nodes and
five-significant-digit dimensionless eigenvalues
\begin{equation}
\boxed{K_1=3.0000,\quad K_2=5.0000,\quad K_3=7.0000}.
\end{equation}
Since $E=(\hbar\omega/2)K$,
\begin{equation}
\boxed{E_1=1.5000\,\hbar\omega,\quad
E_2=2.5000\,\hbar\omega,\quad
E_3=3.5000\,\hbar\omega}.
\end{equation}
Alternating parity and rejecting growing tails
distinguishes the successive bound states.

\problem{2.57}{Numerical shooting for four infinite-well energies}
Scale the box by $y=x/a\in[0,1]$ and set
$\varepsilon=2ma^2E/\hbar^2$. The equation becomes
\begin{equation}
-u''(y)=\varepsilon u(y),\qquad
u(0)=0,\quad u'(0)=1,\quad u(1)=0.
\end{equation}
With trial $\varepsilon>0$,
$u(y)=\sin(\sqrt{\varepsilon}\,y)/\sqrt\varepsilon$.
The shooting residual is
$u(1)=\sin\sqrt{\varepsilon}/\sqrt\varepsilon$,
whose roots give
\begin{equation}
\boxed{\varepsilon_1=9.8696,\quad
\varepsilon_2=39.478,\quad
\varepsilon_3=88.826,\quad
\varepsilon_4=157.91}.
\end{equation}
The exact benchmark is $\varepsilon_n=n^2\pi^2$.
The corresponding eigenfunctions have $n-1$
interior nodes.

\problem{2.58}{Simple estimate of metallic cohesion}
An isolated electron in a width-$a$ box has
ground energy $e_1=\pi^2\hbar^2/(2ma^2)$.
For $N$ separate atoms with one electron each,
$E_{\rm sep}=Ne_1$.
Place the $N$ electrons (temporarily ignoring spin)
in a single box of width $Na$, one in each
of its lowest $N$ levels. The total is
\begin{align}
E_{\rm metal}
&=\sum_{n=1}^N\frac{n^2\pi^2\hbar^2}{2mN^2a^2}\\
&=e_1\frac{N(N+1)(2N+1)}{6N^2}
=e_1\frac{(N+1)(2N+1)}{6N}.
\end{align}
Thus the energy required to separate the atoms is
$E_{\rm coh}=E_{\rm sep}-E_{\rm metal}$.
For large $N$,
\begin{equation}
\boxed{\frac{E_{\rm coh}}{N}\longrightarrow
\frac23 e_1=\frac{\pi^2\hbar^2}{3ma^2}}.
\end{equation}
At $a=4$ \AA{} with the electron mass,
\begin{equation}
\boxed{E_{\rm coh}/N\simeq1.57\ \mathrm{eV}}.
\end{equation}
This toy model omits spin, three-dimensional
bands and ion-core structure, but captures
the lowering of kinetic energy when electrons
delocalize across a larger volume.

\problem{2.59}{Quantum bouncing ball: Airy eigenstates}
For a particle above a reflecting floor, let
$V=mgx$ when $x>0$ and infinite below.
Define
\begin{equation}
a_g=\left(\frac{2m^2g}{\hbar^2}\right)^{1/3},
\quad z=a_gx,\quad
\lambda=\frac{a_gE}{mg}
=\frac{2mE}{\hbar^2a_g^2}.
\end{equation}
The Schrödinger equation becomes
$-y''(z)+zy(z)=\lambda y(z)$,
whose general solution is
$C\operatorname{Ai}(z-\lambda)+
D\operatorname{Bi}(z-\lambda)$.
The $\operatorname{Bi}$ term diverges as
$z\to+\infty$, so $D=0$. The floor enforces
$\operatorname{Ai}(-\lambda_n)=0$.
The first and tenth positive Airy-zero magnitudes
are
\begin{equation}
\boxed{\lambda_1=2.338107410,\qquad
\lambda_{10}=12.828776753}.
\end{equation}
The Airy normalization identity gives
\begin{equation}
\int_0^\infty\operatorname{Ai}(z-\lambda_n)^2\dd z
=\operatorname{Ai}'(-\lambda_n)^2,
\qquad
\boxed{\psi_n(x)=
\frac{\sqrt{a_g}\operatorname{Ai}(a_gx-\lambda_n)}
{|\operatorname{Ai}'(-\lambda_n)|}}.
\end{equation}
In units $a_g=1$, the normalizing
prefactors are approximately $1.426105$
and $0.936510$ for the two states.
Their orthogonality follows from the
self-adjoint Sturm--Liouville equation,
and can also be verified numerically.
Virial and Airy integral identities give
\begin{equation}
\avg z=\frac23\lambda,\qquad
\avg{z^2}=\frac8{15}\lambda^2,\qquad
\avg{p^2}=\hbar^2a_g^2\,\frac{\lambda}{3},
\quad\avg p=0.
\end{equation}
Therefore
\begin{equation}
\boxed{\sigma_x=\frac{2\lambda}{\sqrt{45}\,a_g},
\quad \sigma_p=\hbar a_g\sqrt{\frac{\lambda}{3}},
\quad \frac{\sigma_x\sigma_p}{\hbar}
=\frac{2\lambda^{3/2}}{\sqrt{135}}}.
\end{equation}
For $\lambda_1$ this is $0.61540$, and
for $\lambda_{10}$ it is $7.90935$,
both above $1/2$.
Classically, the particle has a height cutoff
$z=\lambda$ and residence-time density
\begin{equation}
\boxed{\rho_C(z)=\frac1{2\sqrt{\lambda(\lambda-z)}}},
\quad0<z<\lambda.
\end{equation}
This diverges integrably at the classical
turning point, while the quantum density
$|\psi_n|^2$ oscillates and extends
beyond it; the coarse-grained high-$n$
profile approaches the classical envelope.
\check{The companion Python numerics script
plots these comparisons in the scaled coordinates.}

\problem{2.60}{Inverse-square attraction and the boundary-condition anomaly}
For $V=-\alpha/x^2$ on $x>0$, the constants
$m,\hbar,\alpha$ provide no independent length
scale: $\beta=2m\alpha/\hbar^2$ is dimensionless.
Hence there is no nonzero energy uniquely fixed
by dimensional analysis alone.
For $E=-\hbar^2\kappa^2/(2m)$,
\begin{equation}
\psi''+\frac{\beta}{x^2}\psi=\kappa^2\psi.
\end{equation}
If $\psi(x)$ solves this for $\kappa$,
$\psi(sx)$ solves it for $s\kappa$, giving
the formal scaling $E\mapsto s^2E$.
Set $g=\sqrt{\beta-1/4}$, assumed real
($\beta>1/4$). Let $\psi(x)=\sqrt x\,f(\kappa x)$;
the equation for $f(z)$ is
\begin{equation}
z^2f''+zf'-(z^2-g^2)f=0.
\end{equation}
The decaying solution is the modified Bessel
function $K_{\ii g}$:
\begin{equation}
\boxed{\psi_\kappa(x)=C\sqrt x\,K_{\ii g}(\kappa x)}.
\end{equation}
Its norm uses the identity
$\int_0^\infty z\,K_{\ii g}(z)^2\dd z
=\pi g/(2\sinh\pi g)$, so
\begin{equation}
\boxed{C=\kappa\sqrt{\frac{2\sinh(\pi g)}{\pi g}}}.
\end{equation}
As $x\to0$, the wavefunction behaves as
$\sqrt x$ times an oscillatory function of $\ln x$,
with infinitely many nodes approaching zero.
\textit{Important mathematical qualification:}
square integrability alone does not choose
a self-adjoint Hamiltonian at the singular origin.
A fixed self-adjoint extension supplies a
boundary phase/length scale, so formal
continuous rescaling need not preserve its
domain. This is the scale anomaly of the
attractive inverse-square problem.
If an infinite wall is instead placed at
$x=\ell>0$, the new condition is
\begin{equation}
\boxed{K_{\ii g}(\kappa\ell)=0}.
\end{equation}
Its largest positive $\kappa$ root
gives the most negative bound level.
Writing $z_*(g)=\kappa_0\ell$,
\begin{equation}
\boxed{E_0=-\frac{\alpha}{\ell^2}
f(\beta)},\qquad
\boxed{f(\beta)=\frac{z_*(\sqrt{\beta-1/4})^2}{\beta}}.
\end{equation}
This displays explicitly how a length parameter
restores a well-defined energy scale.

\problem{2.61}{Finite-difference Hamiltonian for a box}
Discretize the interval $[0,a]$ into $N+1$
equal segments of length $\Delta x=a/(N+1)$.
With Dirichlet endpoints $\psi_0=\psi_{N+1}=0$,
the $N\times N$ tridiagonal matrix is
\begin{equation}
\boxed{H_N=\frac{\hbar^2}{2m(\Delta x)^2}
\begin{pmatrix}
2&-1&0&\cdots\\
-1&2&-1&\cdots\\
0&-1&2&\ddots\\
\vdots&\vdots&\ddots&\ddots
\end{pmatrix}}.
\end{equation}
In units $E_*= \hbar^2/(2ma^2)$,
the dimensionless eigenvalues and eigenvectors
are exactly
\begin{equation}
\boxed{\varepsilon_n(N)=4(N+1)^2
\sin^2\left(\frac{n\pi}{2(N+1)}\right)},\quad
\boxed{v^{(n)}_j=
\sqrt{\frac2{N+1}}\sin\frac{jn\pi}{N+1}},
\end{equation}
for $n,j=1,\ldots,N$.
This follows by applying the discrete second
difference to a sine eigenvector and using
$\sin(\theta+\delta)+\sin(\theta-\delta)
=2\sin\theta\cos\delta$.
For small sizes:
\begin{equation}
\begin{array}{c|l}
N & \{\varepsilon_n\}\\\hline
1& 8\\
2&9,\ 27\\
3&9.372583,\ 32,\ 54.627417
\end{array}
\end{equation}
For $N=10$ and $100$, the lowest five are:
\begin{center}
\begin{tabular}{rrrr}
\toprule
$n$ & $\varepsilon_n(10)$&$\varepsilon_n(100)$
& $\varepsilon_n(\infty)=n^2\pi^2$\\
\midrule
1&9.802700&9.868809&9.869604\\
2&38.416645&39.465687&39.478418\\
3&83.523702&88.762003&88.826440\\
4&141.469567&157.710064&157.913670\\
5&207.559809&246.243168&246.740110\\
\bottomrule
\end{tabular}
\end{center}
The convergence is of order $(\Delta x)^2$
for fixed low-lying $n$, and the eigenvectors
approach the corresponding continuum sine
functions.

\problem{2.62}{Numerical energies for a sinusoidal well bottom}
Let the infinite well occupy $0<x<a$
with $V(x)=500E_*\sin(\pi x/a)$,
$E_*=\hbar^2/(2ma^2)$.
For $N=100$ interior grid points
$x_j=ja/(N+1)$, use the dimensionless
tridiagonal matrix
\begin{equation}
\boxed{H_{jj}/E_*=2(N+1)^2+
500\sin\frac{\pi j}{N+1},\quad
H_{j,j\pm1}/E_*=-(N+1)^2}.
\end{equation}
Diagonalization yields the three lowest levels
(in units of $E_*$):
\begin{equation}
\boxed{\varepsilon_1\simeq304.65956,\quad
\varepsilon_2\simeq304.80037,\quad
\varepsilon_3\simeq476.16270}
\quad(N=100).
\end{equation}
The numerical eigenvectors are normalized by
$\sum_j|v_j|^2=1$ and can be converted to
continuum-normalized wavefunctions with
$\psi_j=v_j/\sqrt{\Delta x}$.
The first two states are close in energy
because the strong central positive hump
creates two effective edge wells with a
small tunneling splitting.
\check{The indicated $N=100$ matrix results
are specific discretization outputs, not
continuum-exact eigenvalues; the companion
script recomputes them and generates plots.}

\problem{2.63}{Oscillator partition function and Einstein heat capacity}
For equilibrium probabilities
$P_n=\ee^{-\beta E_n}/Z$ with
$Z=\sum_n\ee^{-\beta E_n}$,
\begin{equation}
-\frac{\partial}{\partial\beta}\ln Z
=-\frac1Z\frac{\partial Z}{\partial\beta}
=\frac1Z\sum_n E_n\ee^{-\beta E_n}
=\boxed{\bar E}.
\end{equation}
With oscillator energies $E_n=(n+1/2)\hbar\omega$,
the geometric series gives
\begin{equation}
\boxed{Z=\frac{\ee^{-\beta\hbar\omega/2}}
{1-\ee^{-\beta\hbar\omega}}}.
\end{equation}
Differentiating its logarithm yields
\begin{equation}
\boxed{\bar E=\frac{\hbar\omega}{2}
+\frac{\hbar\omega}{\ee^{\beta\hbar\omega}-1}
=\frac{\hbar\omega}{2}
\frac{1+\ee^{-\beta\hbar\omega}}
{1-\ee^{-\beta\hbar\omega}}}.
\end{equation}
For three oscillators per atom, define
$\theta_E=\hbar\omega/k_B$.
Then $C=3\,\partial\bar E/\partial T$ gives
\begin{equation}
\boxed{C=3k_B
\left(\frac{\theta_E}{T}\right)^2
\frac{\ee^{\theta_E/T}}
{(\ee^{\theta_E/T}-1)^2}}.
\end{equation}
At $T\gg\theta_E$,
$\ee^{\theta_E/T}-1\sim\theta_E/T$, hence
$C\to3k_B$ (classical Dulong--Petit limit).
At $T\ll\theta_E$,
$C\sim3k_B(\theta_E/T)^2\ee^{-\theta_E/T}\to0$.
The textbook's diamond curve illustrates
this temperature dependence qualitatively;
the model does not capture all real-crystal
phonon dispersions.

\problem{2.64}{Power-series solution of Legendre's equation}
For real $\ell\ge0$, start with
\begin{equation}
(1-x^2)f''-2xf'+\ell(\ell+1)f=0,\qquad
f(x)=\sum_{n=0}^{\infty}a_nx^n.
\end{equation}
Substituting
$f''=\sum_n(n+2)(n+1)a_{n+2}x^n$
and collecting the coefficient of $x^n$,
\begin{equation}
(n+2)(n+1)a_{n+2}
+[\ell(\ell+1)-n(n+1)]a_n=0.
\end{equation}
Therefore
\begin{equation}
\boxed{a_{n+2}=
\frac{n(n+1)-\ell(\ell+1)}
{(n+1)(n+2)}\,a_n}.
\end{equation}
Even and odd coefficients form independent
series. For integer $\ell$, the series
of matching parity truncates at $n=\ell$.
The resulting low-degree polynomial solutions,
normalized by $P_\ell(1)=1$, are
\begin{align}
\boxed{P_0(x)}&=\boxed{1},&
\boxed{P_1(x)}&=\boxed{x},\\
\boxed{P_2(x)}&=\boxed{\frac12(3x^2-1)},&
\boxed{P_3(x)}&=\boxed{\frac12(5x^3-3x)}.
\end{align}
For noninteger $\ell$, an isolated parity
series does not terminate and generically has
a logarithmic endpoint singularity. More general
linear combinations, however, include the
non-polynomial Legendre function $P_\ell(x)$,
which is regular at one endpoint. Imposing
regularity at \emph{both} endpoints in the usual
self-adjoint eigenvalue problem selects
nonnegative integer $\ell$.
